\section{Lean formalization: definitions, proof targets, and remaining work}
\label{sec:formalization}
\subsection{The object-theory endpoint}
The intended endpoint is a theorem whose type is an actual derivation in
the arithmetic calculus:
\[
       (\varnothing:\operatorname{List}(\operatorname{Fm}\ \operatorname{PTm}\ 0))
               \ \vdash_{\operatorname{PRAAx}}\ \phi.
\]
The repository notation is \texttt{[] $\vdash$[PRAAx] $\phi$}.
It is materially different from proving a theorem in Lean about
\texttt{Real}, \texttt{EuclideanSpace}, or a host-level real-closed field.
The latter may be useful as a metatheoretic check, but it does not supply
an iPRA derivation. In this package the new algebraic candidates already
have the object-theory conclusion type, while the complete scalar-code
backend is explicitly unfinished.

The pinned dependency is ProveIt commit
\nolinkurl{bbb85238dac533fc4b7b84a750cca864fe35fed5}, with Lean
\texttt{leanprover/lean4:v4.32.0} and Mathlib commit
\nolinkurl{81a5d257c8e410db227a6665ed08f64fea08e997}.
The exact upstream paths and blob hashes that were inspected are recorded
in \nolinkurl{sources/provenance.json}. Its source files, rather than a
memory of an older API, determined the imports and tactic calls below.

\subsection{Literal terms and predicates supplied in the archive}
The file \texttt{lean/EuclideanIPRA/GeometrySyntax.lean} defines
\texttt{AlgebraicSymbols}. Its fields have types \texttt{PR k}, actual
primitive-recursive program syntax of the repository, not arbitrary
functions supplied with no effectivity proof. The fields are validity,
equality, strict order, zero and one codes, addition, negation,
multiplication, reciprocal, and nonnegative square root. A point over
$n$ free variables is a pair of \texttt{PTm n} terms. The file defines:
\begin{itemize}
\item point validity and equality, vector addition and subtraction, scale,
quarter-turn, dot product, determinant, squared norm, and squared distance;
\item collinearity, positive orientation, noncollinearity, distinctness,
betweenness, right angles, parallel directions, and incidence with a line
or a circle;
\item length, cosine, squared sine, and angle equality, with nonzero-ray
hypotheses required when using these raw definitions;
\item explicit terms for midpoint, centroid, perpendicular foot,
reflection, line intersection, circumcenter, orthocenter, incenter,
the two line--circle candidates, and inversion;
\item concrete formula definitions, including \texttt{lineMeetSpec} and
\texttt{circumcenterSpec}, along with \texttt{thalesSpec},
\texttt{lineCircleSpec}, and \texttt{ptolemyInequalitySpec}.
\end{itemize}
These are definitions and targets. No instance of the complete
\texttt{AlgebraicSymbols} record is provided, and no field or root laws
are inserted as unproved axioms. The record isolates the implementation
task precisely: reify the bounded algorithms of
Section~\ref{sec:foundation} and prove their advertised contracts in
\texttt{Prf PRAAx}.

The raw specifications deliberately expose only parts of the eventual
theorem contracts. \texttt{lineMeetSpec} and \texttt{circumcenterSpec}
state incidence or equidistance; validity and uniqueness remain separate
obligations. \texttt{lineCircleSpec} asserts membership of the two
candidate terms when the discriminant is nonnegative; completeness,
the negative case, and the one-versus-two count must be added.
\texttt{thalesSpec} is the squared-distance equivalence, including
endpoints; its geometric right-angle export requires distinct endpoints
and nonzero rays. Finally \texttt{angleEq} compares cosines only.
Its equivalence with the report's interior angle-pair equality needs
the nonnegative sine branch and the proved scalar laws. None of these
raw formulas is presented as an already proved complete geometric theorem.

For example, the printed circle chapter uses a squared radius $\rho$;
the Lean predicate \texttt{onCircle p o r} takes a radius $r$ and compares
the squared distance with $r^2$. They agree after substituting
$\rho=r^2$ and proving $r>0$. This simple representation conversion
must be made explicitly when exporting a theorem. Similarly the report's
comparison-zero equality is the Boolean program \texttt{equal} returning
one in the syntax file. Literal equality of scalar record numbers is not
used as geometric equality.

\subsection{The proposed arithmetic derivations}
\texttt{lean/EuclideanIPRA/Derivations.lean} has 42 theorem declarations:
40 generic integer-coordinate lemmas and two universally closed versions.
Their proofs call the existing \texttt{ipra\_ring} and
\texttt{ipra\_zle} infrastructure. The latter turns a nonnegative slack
certificate into an object-theory inequality. For example, the
Cauchy--Schwarz slack is the square of the determinant. In mathematical
notation the two generic targets are
\[
 \Gamma\vdash N(u)N(v)=(u\cdot v)^2+[u,v]^2,
 \qquad \Gamma\vdash (u\cdot v)^2\le N(u)N(v).
\]
The actual source patterns are:
\begin{lstlisting}
theorem lagrange (u v : Vec2 n) :
    Gamma |- [PRAAx] ZT.eq (sqnorm u * sqnorm v)
      (ZT.sq (dot u v) + ZT.sq (det u v)) := by
  unfold sqnorm dot det
  ipra_ring

-- In the actual source Gamma and |- use Lean's Unicode notation.
-- cauchy_schwarz_squared uses the exact nonnegative certificate:
--   (ZT.toZ Gamma (u.1*v.2-u.2*v.1)) ^ 2
\end{lstlisting}
The listing is a typographic ASCII rendering; the archive contains the
literal Lean source with its correct notation. The closed versions use
eight natural variables for the four signed coordinates and
\texttt{closeAll 8}. Other declarations cover the polarization identity,
Heron, median and Stewart formulas, affine intersection components,
reflection, triangle centers, Thales, circle powers, a Ptolemy algebraic
identity, Varignon, and three-dimensional determinant identities.

\paragraph{Verification boundary.}
Lean and Lake executables and compiled dependencies were unavailable in
this execution environment. \emph{None of these new Lean files was
compiled here.} They are candidate proof-producing source, not reported
kernel-checked theorems. Their absence of \texttt{sorry}, new axioms, and
\texttt{native\_decide} is useful source hygiene, not a substitute for a
successful build. \texttt{lean/build\_status.json} records this explicitly.

\subsection{A concrete order for completing the formalization}
\begin{enumerate}
\item \textbf{Compile the finite identity module.} In an existing checkout
at the pinned commit, build its \texttt{Core.IntCertificates} and
\texttt{Core.Definitions} prerequisites. Run the supplied
\texttt{build\_in\_proveit.py}. It checks the commit and relevant source
cleanliness, then compiles only this packet's modules into a separate
output directory. It neither edits nor fetches upstream sources.
Fix any elaboration errors before advancing.
\item \textbf{Reify finite rational and polynomial algorithms.} Use the
repository's signed-integer and bounded-search terms, representing
lists by its list-coding operations or the explicit encoding in this
report. The obligations are coefficient equality after convolution,
division-with-remainder invariants, Bezout identities, determinant
cofactor identities, and the displayed iteration bounds.
\item \textbf{Implement the scalar readers and comparisons.} Prove
\texttt{Valid} preservation, rational interval nesting, the width and
residual estimates, and the finite sign gap. The difficult equivalence is
comparison-zero versus agreement at every precision. Prove it with an
arbitrary precision free, using the decidability of its rational interval
test, then generalize. Do not invoke a universal evaluator for arbitrary
PR-program codes or induction on that universally quantified agreement.
\item \textbf{Derive field and root laws internally.} Prove agreement is
an equivalence and a congruence, arithmetic is representation-independent,
positive denominator clearing is valid, nonzero inverses satisfy their
equations, and the nonnegative root is both existent and unique.
The annihilator, squarefree, and bounded-grid arguments give the exact
remaining numerical lemmas. Instantiate \texttt{AlgebraicSymbols} only
after those programs are specified.
\item \textbf{Export elementary geometry.} For each construction prove
validity of the output record along with its equations. Start with line
intersection and perpendicular projection, then circle intersections and
triangle centers. For inequalities first prove the cleared square or
slack identity and then justify every sign-sensitive cancellation.
For angles retain the upper-semicircle branch data. For Ptolemy retain
the four boundary-order signs as well as concyclicity.
\item \textbf{Add variable-size and precision families.} Polygon sums,
triangulations, construction traces, field-degree certificates, and dyadic
volume bounds use QF invariants on their finite lists. Parameterize by
the input list and requested output precision. A theorem about each fixed
polygon size is weaker than the uniform bounded-loop theorem; preserve
that distinction in statement names.
\item \textbf{Close and audit.} Universally close the complete implication
from validity, positivity, orientation, and distinctness hypotheses to
the conclusion. Audit the result type and axiom dependencies. Export a
result as finished only after both the kernel check and the representation
contract have succeeded.
\end{enumerate}

\subsection{What the executed checks establish}
The integer polynomial checker stores a polynomial as an exponent-tuple
dictionary with exact integer coefficients. Addition merges coefficients;
multiplication performs finite convolution; powers iterate that operation.
By distributivity and induction on these finite operations, the resulting
dictionary evaluates to the original expression under every assignment in
every commutative ring. Equal dictionaries therefore certify universal
identities, rather than checking finitely many example points.
The normalization code itself is not formally verified in this package.
The full coefficient records are retained for independent reproduction.

It checked 41 scalar components of 37 named relations. A separate SymPy
script checked 17 circle identity groups, including inscribed-angle
doubling, external tangents, both circle-intersection formulas, inversion,
the Ptolemy orientation calculation, and the butterfly cancellation.
The solid checker uses exact rational pairs $a+b\sqrt5$ to enumerate
supporting faces and validate all five coordinate realizations.
The measurement checker verifies six symbolic identity groups, thirteen
polygon refinement levels with outward rational interval arithmetic, and
512 exact finite-sum and refinement checks. It does not compare any
approximation with floating-point $\pi$.
These checks address algebra and finite examples. They do not alone prove
the sign, representation, root-existence, or complete iPRA-derivability
claims of the human argument. The verification directory records the
commands, results, and that distinction.
